% Abstract draft. Generated by make_paper_text.py -- do not edit numbers by hand.
\begin{abstract}
Multimodal retrieval-augmented generation can answer a document question from optical character
recognition (OCR) text, from a retrieved page image, or from both, at different GPU cost. Serving it at scale is
a GPU provisioning problem: generation compute grows with the request rate, and late-interaction visual indexes grow
with the corpus, to about 246\,GiB at a million pages. We present the Cost-Aware Visual Router (CAVR), which
chooses the answer path per query with a run-time cost weight. On one RTX 4090, we run every path on
8,150 DocVQA and InfographicVQA questions with one shared generator, measuring average normalised Levenshtein
similarity (ANLS), GPU-seconds, energy and latency. Its content-aware variant CAVR-CE, a cross-encoder (CE) over
the retrieved passages, uses
29\% less GPU compute than always using both contexts, at 0.008 lower mean ANLS, and 20\%
less than a quality-blind visual/hybrid mix at equal ANLS. In a load test simulated from
measured service times, it cuts the GPUs needed for 100 queries/s at a p95 latency of 2\,s from 106 to
72; always using the page image alone needs 64 at lower ANLS. On 1,499
held-out questions, the frozen router still saved 35\% against always using both contexts, but the page image
alone came within 0.005 ANLS at 16\% less compute. Keyword-derived route labels agree with
measured outcomes no better than chance. On two further GPU nodes, index sharding
scales near-linearly, and two GPUs batch-search an index of a million pages in host memory at near-resident
cost per query, though not interactively.
\end{abstract}
